\title{Parameter-Efficient Orthogonal Finetuning via Factorization}

\begin{abstract}
Large foundation models are becoming increasingly prevalent, but training them from the ground up is prohibitively costly. Consequently, adapting these robust models efficiently to specific downstream tasks is gaining importance. In this work, we investigate a principled finetuning approach—Orthogonal Finetuning (OFT)—for downstream task adaptation. Although OFT exhibits strong generalization capabilities, it still requires a substantial number of trainable parameters due to the large dimensionality of orthogonal matrices. To mitigate this, we begin by analyzing OFT through the lens of information transmission, identifying several crucial desiderata that foster improved parameter efficiency. Drawing inspiration from the Cooley-Tukey fast Fourier transform algorithm's ability to facilitate efficient information transfer, we introduce a compact orthogonal parameterization leveraging butterfly structures. We integrate this parameterization into OFT, leading to a novel parameter-efficient finetuning method termed Orthogonal Butterfly (BOFT). By encompassing OFT as a particular instance, BOFT establishes a broader orthogonal finetuning framework. Finally, we perform an extensive empirical evaluation by adapting large vision transformers, large language models, and text-to-image diffusion models to a variety of downstream vision and language tasks.
\end{abstract}

\section{Introduction}

Recent models such as ChatGPT~\cite{brown2020language,chen2021evaluating} and Stable Diffusion~\cite{rombach2022high} showcase the impressive generalization capabilities of large foundation models. The rapid advancements in these models' performance come alongside a substantial increase in parameter counts (e.g., GPT-3 contains roughly 175 billion parameters). As a consequence, training a foundation model from scratch has become increasingly difficult for researchers. Thus, further progress in the field depends on the capacity to adapt such models without full retraining. In other words, efficiently customizing pre-existing foundation models for downstream tasks is essential. There exist mainly three strategies for efficient task adaptation: \emph{model finetuning}~\cite{hulora2022,zaken2022bitfit,guo2020parameter,raffel2020exploring,qiu2023controlling,zhang2022adaptive,chavan2023one}, which keeps the model architecture fixed and finetunes a subset of its parameters; \emph{adapter tuning}~\cite{rebuffi2017learning,houlsby2019parameter,pfeiffer2020adapterfusion,lin2020exploring,he2021towards}, involving the addition of trainable parameters appended to the original model; and \emph{prompt tuning}~\cite{li2021prefix,lester2021power}, which appends trainable prefix tokens to the input. Among these techniques, model finetuning stands out as a straightforward yet effective method, and crucially, it introduces no extra inference latency.

The core idea behind model finetuning is to maintain similarity between the pretrained model and the finetuned model according to specific criteria, thereby retaining the knowledge acquired during pretraining. For example, many existing finetuning techniques utilize a small learning rate, which heuristically encourages minimal deviation between the pretrained and finetuned models. Considering the structured characteristics of weight matrices, a more systematic strategy aims to preserve the relational information within these matrices, specifically the pairwise angles between neurons. This concept has led to the development of a novel finetuning framework called Orthogonal Finetuning~(OFT)~\cite{qiu2023controlling}, where neurons within the same layer are transformed using a shared orthogonal matrix, ensuring that the pairwise angles between neurons are provably maintained throughout the finetuning process. Although OFT has shown promising results in terms of generalizability and convergence for finetuning text-to-image diffusion models~\cite{qiu2023controlling}, the large dimensionality of orthogonal matrices results in a substantial number of trainable parameters. To mitigate this, OFT adopts a block-diagonal structure to decrease the parameter count. However, this parameter efficiency comes with trade-offs: the orthogonal matrix possesses a fixed sparsity pattern, and the orthogonal transformations are applied independently within separate blocks. This block-diagonal sparsity, while saving parameters, may introduce unwanted inductive biases (\emph{e.g.}, reducing expressiveness and limiting the ability to approximate classical linear transformations), and crucially, the method for identifying an optimal sparsity pattern remains unclear.

The central challenge in tackling this issue is to produce a dense orthogonal matrix that remains efficient in terms of parameters. Although it might initially appear impossible because a dense orthogonal matrix of dimension $d$ necessitates $\mathcal{O}(d^2)$ parameters, we adopt an innovative strategy by composing a dense orthogonal matrix from multiple factorized sparse matrices. This method is motivated by the insight that reducing the number of trainable parameters can be accomplished by exchanging computation time for memory usage. Since the orthogonal matrix is represented as a product of sparse matrices, these multiplications need to be repeatedly performed during each training iteration. To better understand the matrix factorization, we conceptualize the creation of a dense orthogonal matrix as an information transmission challenge. Within this framework, constructing a dense orthogonal matrix through multiplication of sparse matrices can be interpreted as the process of transmitting information across a grid-structured graph. This perspective on information transmission allows us to devise various strategies for sparse matrix factorization that limit the count of trainable parameters while still being sufficiently expressive to generate dense matrices.

To ensure parameter efficiency within our information transmission framework, we take inspiration from the butterfly structures found in the Cooley-Tukey fast Fourier transform algorithm~\cite{cooley1965algorithm}, where information is exchanged efficiently among different nodes~\cite{klasing1994broadcasting}. Specifically, the butterfly graph used in the Cooley-Tukey algorithm naturally induces a form of sparse matrix factorization known as butterfly factorization. Employing butterfly factorization enables us to construct a $d$-dimensional dense matrix as a product of $\mathcal{O}(\log d)$ sparse matrices, each containing $\mathcal{O}(d)$ non-zero elements. By ensuring each sparse matrix is orthogonal, we achieve an efficient orthogonal parameterization with only $\mathcal{O}(d \log d)$ parameters, which represents a substantial reduction compared to the original parameterization. Leveraging this efficient orthogonal parameterization, we introduce a novel parameter-efficient finetuning technique called Orthogonal Butterfly~(BOFT). BOFT generalizes OFT as a special instance and offers a comprehensive orthogonal finetuning framework. Both the block-diagonal and butterfly structures share a common trait—sparsity. Each introduces sparsity into orthogonal matrices to reduce the effective number of trainable parameters. It is also noteworthy to compare our approach with the low-rank structure in LoRA~\cite{hulora2022}; both low-rank and sparse matrices are principal families of structured matrices~\cite{candes2011robust} that enable parameter efficiency.

In contrast to the block-diagonal structure employed by OFT to balance expressiveness and regularity, BOFT utilizes the butterfly structure to create a more gradual interpolation between matrices from the full orthogonal group (providing expressivity) and identity matrices (ensuring regularity). This approach allows us to identify a more compact hypothesis class within the orthogonal group tailored for downstream tasks. Given the prevalent application of butterfly structures in many rapid linear transforms, such as the discrete Fourier and discrete cosine transforms, we contend that our structured strategy for parameter efficiency introduces an implicit inductive bias that enhances generalizability and mitigates overfitting. Our reasoning stems from the fact that the butterfly structure can readily recover many classical linear transforms, a feat that the block-diagonal structure in OFT cannot achieve. Our key contributions are summarized below:

\begin{itemize}[leftmargin=*,nosep]
    \item We investigate the challenge of achieving parameter efficiency in orthogonal finetuning through a novel information transmission framework. This framework recasts the design of parameter-efficient dense orthogonal matrices as an information transmission problem over a grid-structured graph.
    \item Drawing inspiration from the butterfly structures in the Cooley-Tukey algorithm, we introduce Orthogonal Butterfly, a parameter-efficient orthogonal finetuning technique developed within the information transmission framework.
    \item We offer several theoretical insights explaining why BOFT can substantially decrease the number of trainable parameters while maintaining strong expressivity and stable training. Owing to the matrix factorization, BOFT also features an intriguing weight interpolation (see Figure~\ref{fig:boft_interpolation}).
    \item For the first time, we extend orthogonal finetuning to a variety of tasks beyond controllable text-to-image generation~\cite{qiu2023controlling}, demonstrating its significant promise as a universal model finetuning approach. To this effect, we apply BOFT across several adaptation tasks spanning computer vision and natural language processing. BOFT consistently outperforms existing state-of-the-art methods by a substantial margin, confirming its superior parameter efficiency and generalization capability.
\end{itemize}

\section{Related Work}

\textbf{Parameter-efficient finetuning (PEFT)}. As foundational models (\emph{e.g.}, \cite{brown2020language,radford2021learning,rombach2022high,kirillov2023segment}) continue to grow larger and more capable, there has been significant interest in finetuning these models for downstream tasks in a manner that is efficient with respect to the number of parameters~\cite{treviso2023efficient,ding2023parameter,lialin2023scaling}. Among the various PEFT approaches~\cite{houlsby2019parameter,aghajanyan2020intrinsic,hulora2022,edalati2022krona,wang2022adamix,gheini2021cross,zaken2022bitfit,guo2020parameter,sung2021training,ansell2022composable,lester2021power,li2021prefix,vu2022spot,he2021towards,mao2021unipelt,karimi2021compacter,liu2022few,sung2022lst,chen2022parameter,jia2022visual,chen2022adaptformer,zhang2022neural,jie2023fact,lian2022scaling,luo2023towards}, reparameterization-based techniques~\cite{aghajanyan2020intrinsic,hulora2022,edalati2022krona} are most closely related to our work. LoRA~\cite{hulora2022} finetunes the pretrained weight matrix by adding a product of two low-rank matrices, demonstrating strong results on natural language processing tasks. Since LoRA fixes the rank of all added matrices, several methods \cite{zhang2022adaptive,valipour2022dylora,zhang2023increlora} dynamically adjust the rank across layers to better allocate the parameter budget. Due to its straightforwardness, this low-rank weight reparameterization has gained widespread popularity \cite{dettmers2023qlora,zi2023delta,chavan2023one}. Motivated by how hyperspherical energy relates to generalization~\cite{liu2018learning,liu2021orthogonal}, \cite{qiu2023controlling} introduces orthogonal finetuning (OFT) as an alternative, effective technique for finetuning text-to-image diffusion models. OFT specifically learns an orthogonal matrix to transform neurons within the same layer, achieving improved generalization and consistently more stable training than LoRA. Although OFT attains strong performance, it typically involves more trainable parameters than LoRA. Hence, enhancing the parameter efficiency of OFT is a valuable objective. Furthermore, it remains uncertain whether OFT can be applied to a broader range of adaptation tasks beyond controlling text-to-image diffusion models~\cite{qiu2023controlling}. BOFT increases OFT’s parameter efficiency by leveraging butterfly factorization. This advancement enables us to showcase the potential of orthogonal finetuning across general adaptation tasks.

\textbf{Butterfly structures}. The radix-2 Cooley-Tukey algorithm decomposes the $N$-point discrete Fourier transform recursively into two $\frac{N}{2}$-point discrete Fourier transforms, resulting in a butterfly structure expressible as a product of several sparse matrices (known as a butterfly matrix). Butterfly matrices have been utilized to parameterize orthogonal matrices to eliminate pivoting in Gaussian elimination and enhance efficiency~\cite{parker1995random}, to stabilize training of recurrent neural networks~\cite{jing2017tunable}, and in kernel approximation~\cite{munkhoeva2018quadrature}. Works such as \cite{dao2019learning,dao2019kaleidoscope} employ butterfly parameterizations to learn fast linear transforms. Additionally, \cite{chen2022pixelated,dao2022monarch} use butterfly matrices to facilitate efficient neural network training. Butterfly structures have also proven useful in fast matrix-vector multiplication~\cite{michielssen1996multilevel,o2010algorithm}, data-sparse matrix approximation~\cite{li2015butterfly}, and network transmission~\cite{klasing1994broadcasting,li2003linear,rajasingh2016transmission}. Differing from prior work, our focus lies in exploiting butterfly structures to improve the parameter efficiency of OFT.

\section{An Information Transmission View on Orthogonal Finetuning}

We begin by outlining some foundational aspects of OFT. To finetune the pretrained weight matrix, OFT redefines the new weight matrix as the product of a learnable orthogonal matrix and the fixed pretrained weight matrix. Unlike LoRA, which updates weights by adding a low-rank matrix, OFT updates weights multiplicatively using an orthogonal matrix. For parameter efficiency, LoRA employs a low-rank structure, whereas the original OFT utilizes a block-diagonal orthogonal structure~\cite{qiu2023controlling} (the smaller the diagonal blocks, the greater the parameter efficiency of OFT). An intuitive comparison is shown in Figure~\ref{fig:comp_lora}. The rationale behind applying an orthogonal transformation to finetune the weight matrix is to maintain the pair-wise angles between neurons~\cite{liu2018learning,liu2021orthogonal,qiu2023controlling}, thereby largely preserving the semantic knowledge acquired during pretraining. Specifically, OFT learns an orthogonal matrix $\bm{R}\in\mathbb{R}^{d\times d}$ for a pretrained linear layer $\bm{W}^0\in\mathbb{R}^{d\times n}$, modifying the forward computation from $\bm{z}=(\bm{W}^0)^\top\bm{x}$ to $\bm{z}=(\bm{R}\bm{W}^0)^\top\bm{x}$, where $\bm{x}\in\mathbb{R}^{d}$ and $\bm{z}\in\mathbb{R}^{n}$ denote the input and output vectors, respectively. To enable zero initialization, OFT initializes $\bm{R}$ as the identity matrix. To maintain the orthogonality of $\bm{R}$ during finetuning, we follow \cite{qiu2023controlling,liu2021orthogonal} and use Cayley parameterization, i.e., $\bm{R}=(\bm{I}+\bm{Q})(\bm{I}-\bm{Q})^{-1}$, where $\bm{Q}$ is a skew-symmetric matrix satisfying $\bm{Q}=-\bm{Q}^\top$. To ensure parameter efficiency, the block-diagonal structure is enforced by expressing the orthogonal matrix $\bm{R}$ as $\text{diag}(\bm{R}_1,\bm{R}_2,\cdots,\bm{R}_r)$, where each $\bm{R}_i\in\mathbb{R}^{b\times b}$ is a small orthogonal matrix and $br=d$. Although this block-diagonal design improves parameter efficiency, it imposes the assumption that the dimensions of a neuron (i.e., a column vector of $\bm{W}^0$) are partitioned into $r$ groups, with different orthogonal transformations applied independently to each group. Despite the empirical success of this block-diagonal approach, dividing neuron dimensions into groups based solely on their indices is arbitrary, which motivates the desirability of dense orthogonal matrices. This leads to a natural question: \emph{Is it possible to construct a dense orthogonal matrix while retaining parameter efficiency?}

To tackle this issue, we suggest constructing a dense orthogonal matrix as a product of several sparse orthogonal matrices. To offer a coherent and intuitive framework for examining the sparsity structure in orthogonal matrix factorization, we reinterpret the task of generating a dense orthogonal matrix in OFT as a problem of information transmission. More precisely, producing a dense matrix $\bm{R}\in\mathbb{R}^{d\times d}$ via the product of $m$ square matrices \(\thickmuskip=2mu \medmuskip=2mu \bm{R}=\bm{B}_m\bm{B}_{m-1}\cdots\bm{B}_1\) can be seen as passing information through a grid of $d\times (m+1)$ nodes, as depicted in Figure~\ref{fig:info_trans}. This information transmission perspective is motivated by the idea that a $d$-dimensional dense square matrix represents a dense connectivity between $d$ nodes and another set of $d$ nodes. For the matrix $\bm{R}$, a zero entry $R_{ij}$ indicates that information cannot flow from the $j$-th node to the $i$-th node, whereas a nonzero $R_{ij}$ signifies successful transmission. Hence, decomposing the dense matrix $\bm{R}$ into a product of matrices $\bm{B}_m\bm{B}_{m-1}\cdots\bm{B}_1$ can also be interpreted as performing a sequence of information exchanges governed by the graphs defined by each $\bm{B}_i$ for $i=1,\ldots,m$. The flow of information begins with $\bm{B}_1$ and proceeds through to $\bm{B}_m$. As a specific illustration in Figure~\ref{fig:info_trans}, we examine the factorization \(\bm{R}=\bm{B}_5\bm{B}_4\bm{B}_3\bm{B}_2\bm{B}_1\) with its sparsity patterns and corresponding graph structure shown. The graph shown in Figure~\ref{fig:info_trans} emerges from unfolding the matrix multiplication. In this induced graph, each matrix $\bm{B}_i$ represents the connectivity from nodes at level $i$ to nodes at level $i+1$. More explicitly, the element $(j_1,j_2)$ of $\bm{B}_i$ indicates whether there is a directed edge from the $j_2$-th node at level $i$ to the $j_1$-th node at level $i+1$ (a zero entry implies no edge). For the product $\bm{B}_5\bm{B}_4\bm{B}_3\bm{B}_2\bm{B}_1$ to yield a dense matrix, every node at the first level must be able to transmit information to all nodes at the sixth level. Restricting attention to $\bm{R}=\bm{B}_4\bm{B}_3\bm{B}_2\bm{B}_1$, which connects source nodes at the first level to receiver nodes at the fifth level, we observe that information from node $1$ cannot reach node $3$, implying a zero entry at position $(3,1)$ in the sparsity pattern of $\bm{B}_4\bm{B}_3\bm{B}_2\bm{B}_1$. Similarly, considering $\bm{R}=\bm{B}_3\bm{B}_2\bm{B}_1$ or $\bm{R}=\bm{B}_2\bm{B}_1$, the induced graph aligns with the corresponding sparsity pattern. In general, for a matrix $\bm{R}\in\mathbb{R}^{d\times d}$ to be dense, the $m$ factor matrices $\bm{B}_m,\ldots,\bm{B}_1$ must represent directed edges on a $d\times (m+1)$ grid where each edge connects nodes only between neighboring levels (i.e., adjacent columns), ensuring that information from every node at level one can be transmitted to every node at level $m+1$.

The information transmission perspective aids in better comprehending the sparsity pattern of factorization matrices in OFT. Figure~\ref{fig:blk_diag} illustrates the block-diagonal configuration of $\bm{R}$ in the original OFT. Although this block-diagonal form reduces the number of trainable parameters, it cannot generate a dense matrix $\bm{R}$. Our objective is to construct a dense orthogonal matrix by combining $m$ sparse orthogonal matrices, while using as few effective trainable parameters as possible. From the information transmission standpoint, the general criteria for achieving this goal are \emph{(i)} \textbf{dense connectivity}: each node in the first layer must have at least one path to every node in the last layer, and \emph{(ii)} \textbf{minimum free edges}: the total number of edges should be minimized subject to the orthogonality constraint. It is important to note that orthogonality imposes a subtle restriction on the edges between consecutive levels. For instance, for each matrix $\bm{B}_i$ to be full-rank (a prerequisite for orthogonality), there must be at least $d$ edges establishing a bijection between all nodes in the $i$-th level and all nodes in the $(i+1)$-th level, setting a lower bound of $d$ edges between adjacent layers (e.g., 4 in the example shown in Figure~\ref{fig:info_trans}). These $d$ edges are essential for orthogonality and should not be included in the count of trainable edges, since these elements are fixed (for example, in a $d \times d$ orthogonal matrix with $d$ non-zero entries, these entries can only be $\pm 1$). Given that orthogonal matrices require fewer parameters than full matrices, the orthogonality constraint induces additional dependencies among edges. As an illustration, for a $2 \times 2$ orthogonal matrix, having a zero in the $(1,1)$ entry necessarily implies a zero at the $(2,2)$ entry as well (i.e., the absence of one edge enforces the absence of another). Consequently, for each feasible edge configuration, orthogonality may enforce the addition or removal of certain edges. Visualizing the non-zero pattern of the composed orthogonal matrix makes the information transmission framework particularly valuable in OFT, since the focus is on the non-zero trainable elements of $\bm{R}$, with their specific values being irrelevant.

A straightforward dense connection between two levels requires $\mathcal{O}(d^2)$ edges (i.e., a single dense orthogonal matrix), resulting in $d^2 - d$ edges (for instance, when $d=4$, there are 12 edges). Figure~\ref{fig:info_trans} illustrates an example of a valid matrix factorization that uses a total of 10 edges, which is actually fewer than what a single dense orthogonal matrix demands. This framework allows us to analyze the parameter efficiency of OFT from a graphical standpoint and makes it simple to identify valid factorizations. We draw inspiration from an intriguing topology known from the Cooley-Tukey algorithm, called butterfly graphs~\cite{cooley1965algorithm}, which can efficiently establish dense connections between $d$ source nodes and $d$ receiver nodes using only $\mathcal{O}(d\log d)$ edges. For example, the structure depicted in Figure~\ref{fig:info_trans} requires 10 edges for dense connectivity, whereas a butterfly network needs just 8 edges. In the following, we describe how the butterfly structure can enhance parameter efficiency.

\section{Orthogonal Parameterization by Butterfly Factorization}

Originally, the butterfly structure is employed in the Cooley-Tukey algorithm to execute the fast Fourier transform. In the Fourier transform, a local modification in the frequency domain can induce a global alteration in the spatial domain, conceptually akin to our problem of information transmission—where every node at the initial level can communicate with all nodes at the final level. The butterfly structure has also gained popularity as a computer network topology~\cite{solihin2015fundamentals,li2011linear} for efficient information exchange. Assuming $k \geq 2$ is a power of two, we define the \emph{butterfly factor} $\bm{B}^F(k)$ as

\begin{equation}\label{eq:but_factor}
\begin{aligned}
    \bm{B}^F(k)=\begin{bmatrix}
    \text{diag}(\bm{d}_1) & \text{diag}(\bm{d}_2) \\
    \text{diag}(\bm{d}_3) & \text{diag}(\bm{d}_4)
    \end{bmatrix} \in \mathbb{R}^{k \times k},
\end{aligned}
\end{equation}

where each $\bm{d}_i \in \mathbb{R}^{\frac{k}{2}}$ for $i = 1, \ldots, 4$ are vectors. For $d = 2^N$, we define the $d$-dimensional \emph{butterfly matrix} $\bm{B}(d) \in \mathbb{R}^{d \times d}$ recursively as

\begin{equation}
\begin{aligned}
    \bm{B}(d) = \tilde{\bm{B}}(d, d) \cdot \begin{bmatrix}
    \bm{B}_1(\frac{d}{2}) & \bm{0} \\
    \bm{0} & \bm{B}_2(\frac{d}{2})
    \end{bmatrix} = \tilde{\bm{B}}(d, d) \tilde{\bm{B}}(d, \frac{d}{2}) \cdots \tilde{\bm{B}}(d, 2),
\end{aligned}
\end{equation}

where $\bm{B}_1(\frac{d}{2})$ and $\bm{B}_2(\frac{d}{2})$ represent two butterfly matrices each of dimension $\frac{d}{2}$. Next, we introduce the \emph{butterfly component} as $\thickmuskip=2mu \medmuskip=2mu \tilde{\bm{B}}(d,k)=\text{diag}(\bm{B}^F_1(k),\cdots,\bm{B}^F_{d/k}(k))$, a block-diagonal matrix sized $d \times d$ with blocks of size $k$, where each $\bm{B}^F_i(k), \forall i$ denotes butterfly factors defined in Equation~\ref{eq:but_factor}. Now, we proceed to utilize the butterfly matrix to parameterize an orthogonal matrix. To do this, it suffices to ensure that every multiplicative factor $\tilde{\bm{B}}(d,k), \forall k$ within the butterfly matrix $\bm{B}(d)$ is orthogonal. We first examine the block-diagonal matrix $\tilde{\bm{B}}(d,2)$ having block size 2, and it is straightforward to ensure orthogonality of $\tilde{\bm{B}}(d,2)$ by using the Cayley transform (or equivalently a 2-dimensional rotation) to parameterize each block, as done in \cite{liu2021orthogonal,qiu2023controlling}. The sparsity pattern of each butterfly component can be interpreted as a permutation of the sparsity pattern of $\tilde{\bm{B}}(d,2)$, thus all butterfly components can be parameterized easily as orthogonal matrices. This approach yields an efficient parameterization of orthogonal matrices composed of multiple $2 \times 2$ orthogonal matrices. Extending butterfly matrices following \cite{chen2022pixelated}, we define a block butterfly component $\tilde{\bm{B}}^b(d,k)$ in which each element in $\bm{d}_i, \forall i$ is replaced by a $b \times b$ matrix. To ensure the block butterfly component $\tilde{\bm{B}}^b(d,2)$ remains orthogonal, we parameterize each $2b \times 2b$ block matrix as orthogonal. The sparsity pattern of the other butterfly components $\tilde{\bm{B}}^b(d,k)$ for $k > 2$ corresponds to a block-wise permutation of the pattern of $\tilde{\bm{B}}^b(d,2)$ and can similarly be converted into orthogonal matrices. Putting all these elements together, the forward pass in BOFT is given by

\begin{equation}
    \begin{aligned}\nonumber
    \bm{z}=\big(\bm{R}(m,b)\cdot\bm{W}^0\big)^\top\bm{x},~~~~\text{s.t.}~\bigg{\{}\bm{R}(m,b)=\prod_{i=1}^m \tilde{\bm{B}}^b_{(d,i)}~~\&~~\underbrace{\big(\tilde{\bm{B}}^b_{(d,j)}\big)^\top\tilde{\bm{B}}^b_{(d,j)}=\tilde{\bm{B}}^b_{(d,j)}\big(\tilde{\bm{B}}^b_{(d,j)}\big)^\top\!=\bm{I}_d}_{\forall j\in [1, m]}\bigg{\}},
    \end{aligned}
\end{equation}

Here, we denote $\tilde{\bm{B}}^b(d,2^{m - i+1})$ as $\tilde{\bm{B}}^b_{(d,i)}$ for convenience, and $\bm{I}_d$ represents the identity matrix of dimension $d$. The orthogonal matrix $\bm{R}(m,b)\in\mathbb{R}^{d\times d}$ is formed by the product of multiple orthogonal butterfly components. For simplicity, we refer to BOFT with $\bm{R}(m,\frac{b}{2})$ as BOFT($m$,$b$), where $b\geq 2$. When $m=1$, BOFT($1$,$b$) simplifies to the block-diagonal OFT~\cite{qiu2023controlling} with block size $b$. BOFT($1$,$d$) corresponds to the original OFT~\cite{qiu2023controlling} featuring an unconstrained full orthogonal matrix. BOFT($\log \frac{2d}{b}$,$b$) employs the block butterfly matrix $\bm{B}^{\frac{b}{2}}(d)$ as $\bm{R}$, resulting in a dense orthogonal matrix $\bm{R}$. Generally, BOFT($m$,$b$) requires $\frac{1}{2}(b-1)dm$ effective trainable parameters for fine-tuning a linear layer of dimensions $d\times n$. If the butterfly matrix is used, i.e., $m=\log d,b=2$, BOFT utilizes $\mathcal{O}(d\log d)$ parameters. In comparison, the original OFT with a full dense orthogonal matrix uses $\mathcal{O}(d^2)$ parameters, and the block-diagonal OFT with $r$ blocks requires $\mathcal{O}(bd)$. Therefore, the original OFT must set the block size as $b=d$ to produce a dense orthogonal matrix, whereas BOFT can achieve this with any $b$.

\textbf{Identity initialization for BOFT}. Fine-tuning techniques generally commence with the exact pretrained model to ensure the fine-tuned model remains close to the original. For instance, LoRA initializes the low-rank weights at zero. In BOFT, all butterfly components are initialized to identity matrices (i.e., the skew-symmetric matrix within the Cayley transform is initialized to zeros).

\textbf{Multiplicative Dropout for BOFT}. LoRA~\cite{hulora2022} additionally introduces a Dropout layer on the low-rank weight update to mitigate overfitting. Conventional Dropout~\cite{srivastava2014dropout} naturally applies to LoRA but not to BOFT because of our multiplicative weight update strategy. To solve this, we propose a multiplicative Dropout tailored for BOFT. Since the orthogonal matrix $\bm{R}(m,b)$ consists of $m$ orthogonal butterfly components, which can be rearranged into $2b\times 2b$ block-diagonal orthogonal matrices, multiplicative Dropout randomly selects $p_1$ percent of these butterfly components and $p_2$ percent of the diagonal blocks within each butterfly component, replacing the selected blocks with identity matrices.

\section{Intriguing Insights and Discussions}

\textbf{Expressivity of BOFT}. The butterfly structure combined with permutations can exactly replicate many well-known fast linear transforms~\cite{dao2019learning,dao2019kaleidoscope} (such as the fast Fourier transform and Hadamard transform); however, the capability of our orthogonal butterfly matrix to approximate a general orthogonal matrix is still unclear. To investigate this, we performed a simulation to approximate a random dense orthogonal matrix~\cite{anderson1987generation} of size $512 \times 512$. The results presented in Figure~\ref{fig:para_eff} represent averages over 10 different random seeds. The y-axis indicates the approximation error, while the x-axis shows the number of effective trainable parameters. Each curve sharing the same color corresponds to BOFT with the same block size, and the leftmost point on each curve represents the error of BOFT($1$,$b$) (i.e., the original block-diagonal OFT with block size $b$). Overall, BOFT achieves greater parameter efficiency compared to OFT. For instance, BOFT($9$,$2$) demonstrates better expressivity than BOFT($1$,$16$) but utilizes considerably fewer parameters. In general, BOFT with smaller $b$ and larger $m$ tends to be more parameter-efficient. For example, BOFT($6$,$4$) requires far fewer parameters yet attains a comparable approximation error to BOFT($2$,$16$). Broadly, the butterfly matrix corresponds to a more structured subset of the orthogonal group (relative to the block-diagonal structure), which theoretically guarantees that BOFT is more expressive than OFT with an identical block size.

\begin{theorem}[Expressivity of BOFT]\label{thm:exp_boft}
    BOFT exhibits greater expressiveness than OFT when the block size is the same. To approximate any orthogonal matrix of dimension $d$, one can multiply butterfly matrices arranged as $\bm{B}_{d-1,1}(d)\bm{B}^\top_{d-1,2}(d)\cdots\bm{B}_{1,1}(d)\bm{B}^\top_{1,2}(d)$, where each $\bm{B}_{i,j}(d)$ is a butterfly matrix for all $i$ and $j$.
\end{theorem}

Theorem~\ref{thm:exp_boft} implies a straightforward extension of BOFT—where the final orthogonal matrix generalizes to $\thickmuskip=2mu \medmuskip=2mu \bm{R}^G(m_1,b_1,m_2,b_2,l)=\bm{R}_{l,1}(m_1,b_1)\bm{R}_{l,2}^T(m_2,b_2)\cdots\bm{R}_{1,1}(m_1,b_1)\bm{R}_{1,2}^T(m_2,b_2)$, with $\bm{R}_{i,j}^T(m,b)$ denoting the orthogonal matrix employed in BOFT. When $m_1=m_2=\log d$, $b_1=b_2=2$, and $l=d-1$, the matrix $\bm{R}^G(m_1,b_1,m_2,b_2,l)$ can represent the full orthogonal group. This kind of matrix composition is also known as the kaleidoscope hierarchy~\cite{dao2019kaleidoscope}. Nevertheless, it should be noted that greater expressivity does not necessarily translate to improved finetuning performance; for example, full finetuning, despite its universal expressivity, frequently results in unsatisfactory outcomes. Balancing expressivity and regularity is crucial for the generalizability of model finetuning. BOFT expands the finetuning parameter space by incorporating structural priors, allowing us to achieve a superior balance between expressivity and regularity.

\textbf{Spectral properties}. Orthogonal finetuning typically achieves superior spectral characteristics compared to LoRA, as it exactly maintains the spectral norm of the pretrained weight matrix $\bm{W}^0$. This can be understood through singular value decomposition: $\bm{W}^0=\bm{U}\bm{\Sigma}\bm{V}^\top$, where $\bm{U},\bm{V}$ are orthogonal matrices and $\bm{\Sigma}$ is a diagonal matrix of singular values. Both OFT and BOFT apply an orthogonal matrix $\bm{R}$ on the left, resulting in the finetuned weights $\bm{R}\bm{U}\bm{\Sigma}\bm{V}^\top$, which leaves the largest singular value (i.e., the spectral norm of $\bm{W}^0$) unchanged. Preserving this property has been demonstrated to significantly enhance training stability and generalization~\cite{miyato2018spectral,yoshida2017spectral}. Additional intriguing mathematical properties are provided in Appendix~\ref{app:sn_property}.

\textbf{Orthogonal finetuning as learning bilinear similarity}. BOFT can be interpreted as learning a bilinear similarity function $\bm{w}^0_i\bm{R}\bm{x}$, where $\bm{w}^0_i$ denotes the $i$-th neuron (i.e., column vector) of the weight matrix $\bm{W}^0$. From this viewpoint, BOFT corresponds to learning a bilinear similarity matrix $\bm{R}$ with a strong constraint (i.e., $\bm{R}$ must be orthogonal), which inherently relates to distance metric learning~\cite{xing2002distance} and bilinear forms~\cite{roman2005advanced}.

\textbf{Inductive bias and generalization in BOFT}. Since the matrix $\bm{R}(m,b)$ in BOFT generally lies within a structured subset of the orthogonal group that restricts the hypothesis space, BOFT naturally introduces an inductive bias. We contend that the structured inductive bias induced by butterfly factorization supports generalization, given that it shares a common structural pattern with many classical linear transforms~\cite{dao2019kaleidoscope}, including the discrete Fourier transform, discrete sine/cosine transform, and Hadamard transform. Furthermore, the sparse matrix factorization employed in BOFT may also introduce an implicit inductive bias~\cite{gunasekar2017implicit,li2018algorithmic,liu2019neural}.

\textbf{Comparison to butterfly-based sparse training}. Several studies~\cite{dao2019learning,dao2019kaleidoscope,chen2022pixelated,dao2022monarch} have explored sparse training using the butterfly parameterization. These works generally emphasize directly reparameterizing weight matrices via the butterfly structure and training neural networks from the ground up. \cite{dao2022monarch} investigates finetuning pretrained weights by initially projecting them onto a variant of butterfly matrices, followed by optimizing these projected components for specific downstream tasks. In contrast, BOFT introduces a distinct finetuning approach that modifies the weights through layer-shared weight matrices.

\input{rebuttal-results}

\section{Concluding Remarks and Limitations}

This paper presents Orthogonal Butterfly (BOFT), a versatile parameter-efficient finetuning technique for foundation models grounded in the butterfly structure. The central idea for enhancing parameter efficiency is to represent a dense orthogonal matrix as a product of multiple sparse orthogonal matrices. To facilitate discovering suitable matrix factorizations, we propose a graphical information transmission framework. Within this framework, we observe that the butterfly structure effectively meets our criteria for sparse orthogonal matrix factorization. We provide empirical evidence of BOFT's effectiveness in finetuning large language models, large vision models, and text-to-image generative models. Our experiments also confirm BOFT’s superiority as a generic finetuning method.

Although BOFT demonstrates strong empirical performance, it is not without limitations. Because the final orthogonal matrix in BOFT is composed by multiplying several orthogonal matrices, the training runtime overhead is somewhat higher compared to OFT. Finding ways to reduce BOFT’s training time remains an open challenge. On the positive side, after finetuning, the BOFT-optimized orthogonal matrices can be directly merged into the pretrained model, resulting in no additional inference latency. Furthermore, it remains uncertain whether the butterfly network represents the most efficient mechanism for information transmission. Our information transmission framework encourages insights from a separate field—computer networking—where the efficiency of network topologies for transmitting information is extensively analyzed. We anticipate that alternative network architectures may offer even more efficient constructions for dense orthogonal matrices.

\end{document}